\documentclass[a4paper]{article}
% Español (México) con XeLaTeX: fontspec para fuentes Unicode, polyglossia
% para la división silábica y los nombres del documento en la variante mexicana.
\usepackage{fontspec}
\usepackage{polyglossia}
\setdefaultlanguage[variant=mexican]{spanish}
% Los nombres chinos van como «pinyin (original)»: xeCJK da a los caracteres
% Han su propia fuente; sin ella XeLaTeX los omite con un aviso.
% FandolSong no cubre algunos caracteres de nombres (U+73FA); WenQuanYi Zen Hei sí.
\usepackage[AutoFallBack=true]{xeCJK}
\setCJKmainfont{FandolSong-Regular.otf}
\setCJKfallbackfamilyfont{\CJKrmdefault}{WenQuanYi Zen Hei}
% polyglossia no traduce los nombres de listings.
\AtBeginDocument{\providecommand{\lstlistingname}{}\renewcommand{\lstlistingname}{Listado}%
  \providecommand{\lstlistlistingname}{}\renewcommand{\lstlistlistingname}{Índice de listados}}
\usepackage{amsmath,amssymb}
\usepackage{graphicx}
\usepackage[margin=2.35cm]{geometry}
\usepackage[most]{tcolorbox}
\usepackage{etoolbox}
\usepackage{listings}
\usepackage{xcolor}
\usepackage{hyperref}
\usepackage{booktabs,longtable,tabularx,array}
\usepackage{subcaption}
\usepackage{float}
\usepackage{enumitem}
\usepackage{microtype}
\usepackage{fancyhdr}
\usepackage{needspace}

\definecolor{kbBack}{HTML}{F2F6FA}
\definecolor{kbFrame}{HTML}{9DB4CC}
\definecolor{kbTitle}{HTML}{52708F}
\definecolor{ibBack}{HTML}{FBF5E7}
\definecolor{ibFrame}{HTML}{C9A961}
\definecolor{ibTitle}{HTML}{7D5F1E}
\definecolor{wbBack}{HTML}{FCF2F0}
\definecolor{wbFrame}{HTML}{D6A096}
\definecolor{wbTitle}{HTML}{9E4F43}
\definecolor{dbBack}{HTML}{F4F7F3}
\definecolor{dbFrame}{HTML}{AEC2AC}
\definecolor{dbTitle}{HTML}{5F7F64}

\tcbset{softbox/.style={
  enhanced,breakable,boxrule=0.8pt,arc=2pt,left=3mm,right=3mm,
  top=2mm,bottom=2mm,coltitle=white,fonttitle=\bfseries,
  toptitle=0.6mm,bottomtitle=0.6mm
}}
\newtcolorbox{knowledgebox}[1]{softbox,colback=kbBack,colframe=kbFrame,colbacktitle=kbTitle,title={#1}}
\newtcolorbox{importantbox}[1]{softbox,colback=ibBack,colframe=ibFrame,colbacktitle=ibTitle,title={#1}}
\newtcolorbox{warningbox}[1]{softbox,colback=wbBack,colframe=wbFrame,colbacktitle=wbTitle,title={#1}}
\newtcolorbox{dialoguebox}[1]{softbox,colback=dbBack,colframe=dbFrame,colbacktitle=dbTitle,title={#1}}

\lstset{
  language=Python,basicstyle=\ttfamily\small,
  keywordstyle=\color{kbTitle}\bfseries,stringstyle=\color{wbTitle},
  commentstyle=\color{dbTitle},breaklines=true,frame=single,numbers=left,
  numberstyle=\tiny\color{gray},captionpos=b,columns=fullflexible,
  keepspaces=true,showstringspaces=false,extendedchars=true
}
\BeforeBeginEnvironment{lstlisting}{\Needspace{12\baselineskip}}

\hypersetup{colorlinks=true,linkcolor=kbTitle,urlcolor=kbTitle,citecolor=wbTitle}
\setlist{itemsep=0.25em,topsep=0.4em}
\setlength{\parindent}{2em}
\setlength{\parskip}{0.25em}
\newcolumntype{L}[1]{>{\raggedright\arraybackslash}p{#1}}

\providecommand{\notetitle}{Notas del curso: [tema del curso]}
\providecommand{\noteauthors}{Elaboradas a partir de los materiales públicos de Stanford CS25}
\providecommand{\notedate}{Versión reescrita en 2026}
\providecommand{\videochannel}{Stanford Online}
\providecommand{\videopublishdate}{2022}
\providecommand{\videoduration}{[duración del video]}
\providecommand{\videourl}{}
\providecommand{\videocoverpath}{cover.jpg}
\providecommand{\coursesite}{https://web.stanford.edu/class/cs25/}
\providecommand{\repourl}{https://github.com/hqhq1025/ai-course-notes}
\providecommand{\skillversion}{wdkns/wdkns-skills@39f1a04}

\newcommand{\figsource}[1]{\par\vspace{-0.3em}{\footnotesize\color{gray}\emph{Fuente: #1}}\par}
\newcommand{\teachervoice}[1]{\begin{knowledgebox}{Nota de clase}#1\end{knowledgebox}}
\newcommand{\term}[1]{\textbf{#1}}

\pagestyle{fancy}
\fancyhf{}
\fancyfoot[C]{\thepage}
\fancyfoot[R]{\footnotesize\color{gray}\href{\repourl}{AI Course Notes}}
\renewcommand{\headrulewidth}{0pt}
\renewcommand{\footrulewidth}{0pt}

\newcommand{\makecscover}{
\begin{titlepage}
\centering
\vspace*{0.7cm}
{\Large Stanford CS25 · Transformers United\par}
\vspace{0.8cm}
{\huge\bfseries \notetitle\par}
\vspace{0.6cm}
{\large \noteauthors\par}
\vspace{0.25cm}
{\large \notedate\par}
\vspace{0.9cm}
\ifdefempty{\videocoverpath}{}{
  \includegraphics[width=0.86\textwidth,height=0.43\textheight,keepaspectratio]{\videocoverpath}\par
}
\vfill
\begin{tcolorbox}[width=0.92\textwidth,colback=black!2!white,colframe=black!45,boxrule=0.7pt,arc=2pt]
\textbf{Autor o canal del video}: \videochannel\par
\textbf{Fecha de publicación}: \videopublishdate\par
\textbf{Duración del video}: \videoduration\par
\textbf{Sitio del curso}: \href{\coursesite}{\nolinkurl{\coursesite}}\par
\ifdefempty{\videourl}{}{\textbf{Enlace del video}: \href{\videourl}{\nolinkurl{\videourl}}\par}
\textbf{Normas de elaboración}: \skillversion\ + las normas source-first / teacher-voice / visual-QA de este repositorio
\end{tcolorbox}
\end{titlepage}
\tableofcontents
\newpage
}

\usepackage{multicol}

\renewcommand{\notetitle}{Lecture 05: De la predicción del siguiente token a la inteligencia artificial de próxima generación — el futuro del preentrenamiento}
\renewcommand{\noteauthors}{Basado en las conferencias oficiales de Shrimai Prabhumoye, restauración visual de 44 videos y subtítulos manuales en inglés organizados}
\renewcommand{\notedate}{CS25 V6 · Apuntes source-first de 2026}
\renewcommand{\videochannel}{Stanford Online}
\renewcommand{\videopublishdate}{Clase del 30 de abril de 2026; subido el 11 de mayo de 2026}
\renewcommand{\videoduration}{57:56 (exposición principal hasta aprox. 49:30, después preguntas y respuestas)}
\renewcommand{\videourl}{https://www.youtube.com/watch?v=e_H_tkpCAK4}
\renewcommand{\videocoverpath}{cover.jpg}

\newcommand{\lecturefigure}[3]{%
\begin{figure}[H]
  \centering
  \includegraphics[width=0.97\textwidth,height=0.48\textheight,keepaspectratio]{slides-images/#1}
  \caption{#2}
  \figsource{#3}
\end{figure}
}

\let\standardtableofcontents\tableofcontents
\renewcommand{\tableofcontents}{%
  \begingroup
  \small
  \setlength{\parskip}{0pt}
  \standardtableofcontents
  \endgroup
}

\begin{document}
\makecscover

\section{El verdadero nuevo eje de escalado: ¿cómo aprender mejor con los mismos datos?}

Esta clase lleva el preentrenamiento desde la pregunta de "cuánto más grande es el modelo y cuántos tokens hay", a la de "\textbf{en qué orden, con qué objetivo y qué tipo de retroalimentación aprende el modelo estos tokens}". El ponente conecta tres líneas de trabajo: \textit{two-phase pretraining} determina la mezcla y el orden de los datos; \textit{front-loading reasoning} decide cuándo entran los datos de razonamiento en el entrenamiento; y RLP (\textit{Reinforcement as a Pretraining Objective}) integra directamente el pensamiento exploratorio en el objetivo del preentrenamiento. Todas ellas modifican el proceso de aprendizaje, no simplemente aumentan la cantidad de corpus.

\subsection{Los cuatro componentes no son cuatro listas independientes}

La clase presenta primero una receta de cuatro componentes para un LLM SOTA. \textit{Smart Data} se encarga de la calidad, cobertura y eliminación de duplicados; \textit{Smart Architecture} del representacional y la estructura computacional; \textit{Smart Algorithms} de los objetivos de optimización, currículo y etapas de entrenamiento; y \textit{Smart Collaboration} hace que el preentrenamiento, el post-entrenamiento, la investigación y la ingeniería dejen de optimizar indicadores locales por separado. Al leer la figura, tenga en cuenta: el enfoque principal de esta clase es el tercer punto, pero cualquier beneficio algorítmico depende de que los otros tres proporcionen datos utilizables, un modelo estable y un sistema ejecutable.

\lecturefigure{slide-002.jpg}{Los cuatro componentes acoplados para construir un LLM SOTA}{Grabación oficial de Stanford, 00:01:21--00:02:09}

\teachervoice{00:01:21--00:02:09, el ponente coloca \textit{smart collaboration} junto a datos, arquitectura y algoritmos. Esta indicación oral es importante: la receta de preentrenamiento no es un solo hiperparámetro en un artículo, sino el resultado de un sistema tras que múltiples equipos compartan definiciones de datos, presupuestos de cómputo y criterios de evaluación.}

\begin{importantbox}{La tesis central de esta clase}
Cuando la escala del modelo y la escala de los datos continúan creciendo, las nuevas capacidades provienen no solo de "más", sino de un \textbf{proceso de aprendizaje más estructurado}: qué datos se observan primero, cuándo se introduce el razonamiento, si se permite que el modelo explore antes de predecir, y si la recompensa puede distinguir entre una exploración efectiva y un monólogo ineficaz.
\end{importantbox}

\subsection{Del \textit{smart data} al \textit{smart algorithms}}

El ponente resume primero su trabajo en datos con una diapositiva, que incluye diálogos matemáticos, extracción matemática de Common Crawl, generación sintética diversa y datos de razonamiento transversales. Aquí no se trata de comentar cada conjunto de datos uno por uno, sino de explicar que los experimentos algorítmicos posteriores no ocurren en el vacío: las etiquetas de calidad de datos, la cobertura temática y la definición de muestras de razonamiento son variables de entrada manipulables por la estrategia de entrenamiento. El estado final de la expansión también indica que estos recursos tienen versiones públicas para facilitar su verificación externa.

\lecturefigure{slide-003.jpg}{Resumen del trabajo de \textit{Smart Data}: matemáticas, extracción web, síntesis diversa y razonamiento transversal}{Grabación oficial de Stanford, 00:02:09--00:03:15}

A continuación, la clase cambia claramente a tres problemas algorítmicos. El primero pregunta "cómo aprovechar al máximo el potencial de los datos existentes", el segundo "si los datos de razonamiento deberían aparecer más temprano", y el tercero "si se puede hacer que el modelo obtenga recompensas mediante exploración durante el preentrenamiento". Las tres líneas de trabajo profundizan capa por capa: primero cambiando la distribución de muestreo, luego la asignación de etapas y finalmente el objetivo del entrenamiento. Si se mezclan en una sola receta, no se puede determinar si los beneficios provienen de los datos, del orden o del \textit{objective}.

\lecturefigure{slide-004.jpg}{Las tres líneas de trabajo algorítmicas de esta clase: dos fases, carga anticipada y RLP}{Grabación oficial de Stanford, 00:03:15--00:03:37}

\begin{knowledgebox}{Términos clave: tres niveles de intervención}
\begin{tabularx}{\textwidth}{L{0.22\textwidth}X X}
\toprule
Intervención & ¿Qué cambia? & ¿Qué se mantiene relativamente comparable? \\
\midrule
\textit{Optimal blend} & Pesos de cada fuente de datos y次数 de repetición & Presupuesto total de tokens y configuración del modelo \\
\textit{Two-phase curriculum} & Orden de aparición de los datos en el progreso del entrenamiento & Pooled conjunto de datos general y función objetivo \\
\textit{Front-loading reasoning} & Etapa en que entran los datos de razonamiento en el flujo de entrenamiento & Presupuesto total fijo de datos de razonamiento y flujo de post-entrenamiento \\
RLP & Si se muestrea \textit{thought} antes de predecir y cómo se calcula la recompensa & Token real siguiente y flujo general de preentrenamiento \\
\bottomrule
\end{tabularx}
\end{knowledgebox}

\subsection{Pascal, Volta, Ampere, Hopper: metáforas pedagógicas de variables controladas}

Cuatro aprendices poseen los mismos datos pero emplean estrategias distintas. Pascal no usa currículo, no utiliza datos ricos en razonamiento y no explora mediante pensamiento; Volta solo hace el currículo; Ampere añade además datos tempranos de razonamiento; y Hopper incorpora finalmente el aprendizaje a través del pensamiento. La matriz final no es una teoría de educación infantil, sino un índice experimental: cada gráfica de resultados posterior compara qué sucede cuando un signo de "x" en una columna cambia por un check.

\lecturefigure{slide-005.jpg}{Estrategias de currículo, carga anticipada y aprendizaje a través del pensamiento correspondientes a los cuatro aprendices}{Grabación oficial de Stanford, 00:03:37--00:06:27}

\teachervoice{00:04:15--00:05:18, el ponente explica la ausencia de currículo con el ejemplo de "leer un artículo histórico al azar, resolver una matemática avanzada y luego leer matemáticas básicas", e ilustra la carga anticipada mencionando cursar AP antes de tiempo. El valor de la metáfora radica en distinguir entre "poseer datos" y "utilizar los datos en la etapa correcta".}

\begin{warningbox}{La metáfora no es evidencia causal}
Las historias sobre Pascal hasta Hopper solo ayudan a memorizar las condiciones experimentales. Los modelos reales no comprenden el significado pedagógico de "fundamentos primero"; si algo es efectivo, hay que volver al presupuesto de datos controlado, el cómputo de entrenamiento, los benchmarks y los intervalos de confianza. Todos los números posteriores describen configuraciones específicas del artículo.
\end{warningbox}

\subsection{Resumen de la sección}

Este capítulo establece un sistema de coordenadas comparativo para toda la clase: \textit{smart data} proporciona fuentes de datos operativas, \textit{smart algorithms} determina los pesos, el orden, las etapas y las recompensas; y Pascal a Hopper organiza tres intervenciones en una serie experimental progresiva. A continuación se aborda primero el paso más conservador: no cambiar el modelo ni la pérdida, sino rediseñar únicamente la mezcla de datos y el currículo.
\section{De «cuántos datos» a «cómo construir la mezcla óptima»}

La escala de los datos sigue aumentando, pero el texto generado por humanos no es infinito. Por ello, el ingeniería de preentrenamiento no puede limitarse a preguntar «¿cuántas más páginas web podemos capturar?», sino que debe interrogar si cada token vale la pena aparecer, cuántas veces debe hacerlo y en qué etapa del entrenamiento: temprana o tardía. La clave aquí es modelar por separado los pesos de mezcla (mixture weights) y el orden temporal; de lo contrario, se mezclan dos fuentes distintas de beneficios.

\subsection{Hay muchas fuentes de datos, pero lo verdaderamente escaso son las oportunidades de entrenamiento efectivas}

Las fuentes como documentos legales, libros, artículos académicos, páginas web y materiales matemáticos presentan diferencias enormes en formato, calidad, escala y grado de repetición. Que un gran corpus tenga muchos tokens no implica que tenga alto valor de entrenamiento; del mismo modo, que un corpus pequeño sea de alta calidad no garantiza que el muestreo infinito evitará el sobreajuste. La primera figura ilustra la diversidad de los datos y la segunda muestra el crecimiento de la escala frente al horizonte de datos generados por humanos; ambas explican por qué se vuelve necesaria la optimización de la receta (recipe optimization).

\lecturefigure{slide-006.jpg}{Composición multifuente de los datos de preentrenamiento para LLMs}{Grabación oficial de Stanford, 00:06:28--00:06:51}

Al interpretar la figura, separe primero «fuentes» de «muestras». Un rastreo web (Web crawl) puede representar la mayoría de los tokens, pero su distribución interna de calidad es muy amplia; por el contrario, los artículos o documentos legales tienen una escala menor, pero pueden proporcionar señales de alta densidad en capacidades específicas. La estrategia de mezcla no consiste en pegar una etiqueta permanente a cada fuente, sino asignar presupuestos combinando filtrado, dominio, dificultad, número de repeticiones y validación downstream.

\lecturefigure{slide-007.jpg}{Crecimiento de la escala de tokens de entrenamiento frente al horizonte de datos generados por humanos}{Grabación oficial de Stanford, 00:06:51--00:07:49}

\begin{knowledgebox}{Lectura de la figura: ¿cómo interpretar la curva de crecimiento azul?}
El eje horizontal representa el tiempo y el vertical la cantidad de datos consumidos para entrenar el modelo. Los puntos históricos y la banda de predicción destacan la transición desde los cientos de billones de tokens de la era de GPT-3 hacia los decenas de billones de trillones. Esto respalda el juicio de recursos según el cual «aumentar simplemente nuevos textos humanos se volverá cada vez más difícil», pero no prueba que un año específico agote exactamente los datos, ni considera el valor de los datos sintéticos, multimodales, trayectorias de interacción y la repetición del entrenamiento.
\end{knowledgebox}

\subsection{Dos preguntas: peso y orden}

Ante el mismo conjunto de fuentes, deben responderse dos problemas ortogonales. Primero, \textbf{cómo ponderar}: con qué probabilidad se extrae cada fuente dentro de un lote (batch). Segundo, \textbf{cómo ordenar}: si esa probabilidad cambia conforme avanza el entrenamiento. El primero constituye la mezcla (blend) y el segundo el currículo (curriculum). Si solo se reporta la proporción total final sin indicar el orden temporal, los lectores externos no podrán replicar la distribución real de entrenamiento.

\lecturefigure{slide-008.jpg}{Dos preguntas para maximizar el potencial de los datos: cómo ponderar y cómo ordenar}{Grabación oficial de Stanford, 00:07:49--00:08:22}

\begin{importantbox}{Diferencia matemática entre Blend y Curriculum}
Sea $D_1,\ldots,D_K$ el conjunto de fuentes de datos. Un blend fijo utiliza pesos constantes $w_k$, cumpliendo $w_k\ge 0$ y $\sum_k w_k=1$; un currículo, en cambio, emplea $w_k(s)$ que varían con el progreso del entrenamiento $s\in[0,1]$. Incluso si la cantidad total de tokens integrada a lo largo de todas las fases es idéntica, diferentes funciones $w_k(s)$ pueden generar modelos distintos.
\end{importantbox}

\subsection{Estimación de calidad, estimación de epochs y creación del blend}

El primer paso para obtener una mezcla óptima es la estimación de calidad: evaluar la efectividad de cada fuente en las capacidades objetivo mediante reglas, clasificadores, rúbricas de LLM o evaluación humana. El segundo paso es la estimación de epochs: determinar cuándo el beneficio marginal se vuelve negativo al repetir cierto dominio 1, 2, 4 o 6 veces. Los pesos finales dependen tanto de la calidad como del número de repeticiones posibles; «alta calidad» no equivale a una licencia de reproducción infinita.

\lecturefigure{slide-009.jpg}{Tres pasos para una mezcla de datos óptima: estimación de calidad, estimación de repeticiones y construcción del blend}{Grabación oficial de Stanford, 00:08:22--00:09:29}

Si $N_k$ es el número único de tokens en la fuente $k$, se planean $e_k$ repeticiones y el presupuesto total es $B$, entonces un peso pedagógico puede expresarse como
\[
w_k=\frac{q_k e_k N_k}{\sum_j q_j e_j N_j},
\qquad
\sum_k w_k=1.
\]
En esta fórmula, $q_k$ denota el coeficiente de calidad empírico, $e_k$ el número estimado de epochs repetibles y $N_k$ la escala de la fuente. Las recetas reales suelen incorporar además un piso de dominio (domain floor), temperatura, deduplicación y límites de muestreo; por tanto, esta ecuación es un marco conceptual y no la única implementación posible en la literatura.

\begin{lstlisting}[caption={Versión pedagógica del pseudocódigo para el abanico de calidad y repeticiones a nivel de dominio}]
for domain in domains:
    quality_score = evaluate_quality_classifier(domain)
    for repeats in [1, 2, 4, 6]:
        candidate = reallocate_budget(
            base_blend, add=(domain, repeats), remove="low_quality_crawl"
        )
        model = train_small_proxy(candidate, fixed_token_budget)
        record(domain, repeats, downstream_score(model))

best_repeats = choose_before_diminishing_returns(records)
final_blend = assemble_domain_level_weights(quality_score, best_repeats)
\end{lstlisting}

\teachervoice{00:25:17--00:26:45, el orador añade en las preguntas y respuestas: el abanico de conteo de repeticiones traslada el presupuesto desde el rastreo de baja calidad hacia el dominio objetivo, lo cual generalmente se hace por categoría de dominio, no ejecutando un conjunto costoso de experimentos para cada fuente de datos independiente.}

\begin{warningbox}{La Calidad no es una propiedad intrínseca del documento}
Un mismo documento puede tener distinto valor según la fluidez del lenguaje, el razonamiento matemático, la generación de código o la búsqueda legal. Clasificadores como FineWeb-Edu solo otorgan un puntaje educativo bajo una rúbrica específica; el ingeniero de entrenamiento debe registrar la rúbrica, la versión del modelo, los umbrales, la composición de dominios y la validación downstream, ya que sin ello «alta calidad» no puede ser revisada.
\end{warningbox}

\subsection{Resumen de la sección}

La botella de cuello en los datos desplaza el foco de la investigación desde «capturar más» hacia «distribuir mejor». Los pesos del blend responden qué porcentaje ocupa cada tipo de dato, mientras que el currículo indica cuándo aparecen; la estimación de calidad y la estimación de epochs miden respectivamente la densidad de señal y el límite superior de repeticiones. El capítulo próximo concretará $w_k(s)$ variable en el tiempo mediante un preentrenamiento de dos fases (two-phase pretraining).


\section{Preentrenamiento en dos fases: primero expandir la cobertura, luego concentrar la calidad}

El preentrenamiento en dos fases no consiste simplemente en dividir el entrenamiento por la mitad, sino en utilizar dos distribuciones de muestreo con objetivos diferentes. La fase 1 expone al modelo a un conocimiento del mundo amplio y diversas expresiones; la fase 2, manteniendo las capacidades generales, aumenta la probabilidad de muestreo de fuentes matemáticas, de código, Wikipedia, de tareas y de otros orígenes de alta calidad. El orden mismo se convierte en una variable optimizable.

\subsection{El mecanismo central del currículo por tramos}

La palabra clave de la fase 1 es la diversidad: grandes cantidades de rastreo web (web crawl), menor repetición de datos de alta calidad y una cobertura de dominios más amplia. La palabra clave de la fase 2 es la calidad: aumentar el número de épocas con datos de alta calidad, utilizando las actualizaciones limitadas tardías para señales de capacidad más densas. Las áreas roja y verde en el gráfico no representan una división moral entre "datos malos" y "datos buenos", sino una asignación de recursos por etapas diferentes.

\lecturefigure{slide-010.jpg}{Preentrenamiento en dos fases: la fase 1 prioriza la diversidad y la fase 2 se centra en la calidad}{Grabación oficial de Stanford, 00:09:29--00:10:14}

El currículo puede escribirse como una distribución de muestreo por tramos:
\[
p(D_k\mid s)=
\begin{cases}
w_k^{(1)}, & 0\le s<s_0,\\
w_k^{(2)}, & s_0\le s\le 1,
\end{cases}
\qquad
w_{\text{HQ}}^{(2)}>w_{\text{HQ}}^{(1)}.
\]
Donde $s$ es el progreso de entrenamiento normalizado, $s_0$ es el límite entre fases (phase boundary), y $w_k^{(1)}$ y $w_k^{(2)}$ son los pesos de los datos para ambas etapas. El informe del artículo indica que utilizar la fase 2 en aproximadamente el último 40\% del entrenamiento produce mejores resultados en su configuración, pero $s_0$ aún debe determinarse mediante experimentos proxy.

\begin{lstlisting}[caption={Muestreador two-phase bajo presupuesto de token fijo}]
def sample_source(step, total_steps, phase_boundary, phase1, phase2):
    progress = step / total_steps
    weights = phase1 if progress < phase_boundary else phase2
    return categorical_sample(weights)

for step in range(total_steps):
    source = sample_source(step, total_steps, 0.60, p1_weights, p2_weights)
    batch = next_batch(source)
    optimize_next_token_loss(batch)
\end{lstlisting}

\subsection{Tres líneas base (baselines) que aíslan tres variables}

La distribución natural muestrea según el número de tokens existentes en cada fuente de datos: las fuentes grandes y de baja calidad ocuparán más actualizaciones debido a su escala. El *optimal blend* con orden aleatorio ya utiliza información de calidad y épocas, pero no altera el orden temporal. El enfoque de dos fases controla simultáneamente la calidad, las épocas y el orden. Solo con un diseño de líneas base así se puede separar el beneficio del "mejor mezcla" del del "mejor currículo".

\lecturefigure{slide-011.jpg}{Distribución natural, mezcla aleatoria óptima y línea base two-phase}{Grabación oficial de Stanford, 00:10:14--00:11:42}

\begin{knowledgebox}{Lectura de la figura: primero observe las tres columnas, luego los nombres del modelo}
Las tres columnas corresponden a Calidad (Quality), Épocas (Epochs) y Orden (Order). La distribución natural no controla ninguno de los tres; el *optimal random blend* controla los dos primeros; *two-phase* controla además el orden. Los nombres Pascal y Volta son solo ayudas mnemotécnicas. Si algún nuevo experimento modifica simultáneamente el tokenizador, la escala del modelo o el número total de tokens, la comparación ya no será limpia según la definición de esta tabla.
\end{knowledgebox}

\subsection{¿Qué responden el 17\% y el 3.4\% respectivamente?}

La sección anterior ya separó las variables de control entre las tres líneas base; esta sección explica con mayor detalle a qué contrafactual corresponde cada uno de los dos números destacados, evitando confundir el "beneficio total de optimizar la mezcla" con el "beneficio aportado únicamente por el orden". En el gráfico de resultados, la mejora relativa promedio de Volta frente a Pascal es de aproximadamente 17\%, donde tanto la mezcla de calidad como el orden cambian. Frente a la línea base que ya posee *optimal blend* pero con orden aleatorio, *two-phase* sigue mostrando una mejora de aproximadamente 3.4\%. El primer número indica la importancia de la receta (*recipe*), mientras que el segundo se acerca más al aporte independiente del orden del currículo.

\lecturefigure{slide-012.jpg}{Resultados de la estrategia two-phase frente a la distribución natural y la mezcla aleatoria óptima}{Grabación oficial de Stanford, 00:11:42--00:12:32}

Si la precisión de la línea base es $a$ y la del nuevo método es $b$, la mejora relativa es
\[
\Delta_{\mathrm{rel}}=\frac{b-a}{a}\times 100\%,
\]
mientras que la mejora absoluta es $\Delta_{\mathrm{abs}}=b-a$ puntos porcentuales. Las menciones de "\% mejor" en el curso se refieren a mejoras relativas y no deben mezclarse con puntos porcentuales absolutos, ni sumarse directamente entre diferentes suites de benchmarks.

\teachervoice{00:27:16--00:27:58, cuando el orador responde a si es factible "primero calidad, luego diversidad", dice que intercambiar las dos fases en un estudio de ablicación (ablation) da resultados peores. Este resultado apoya la dirección actual, pero sigue siendo una experiencia bajo un corpus y régimen de modelo específicos, no una ley universal del aprendizaje.}

\begin{warningbox}{Las dos fases no significan "primero basura, luego joyas"}
La fase 1 aún requiere filtrado y control básico de calidad; la fase 2 tampoco debe limitarse a retener solo datos de puntuación estrecha y alta. Acortar demasiado la cobertura causaría pérdida de alcance; repetir datos de alta calidad demasiado tarde o durante un periodo excesivo produciría rendimientos decrecientes (diminishing returns). En términos de ingeniería, se deben monitorear simultáneamente la cobertura de dominios, la pérdida de validación (*validation loss*), el fenómeno de memorización y la transferencia a tareas posteriores (downstream transfer).
\end{warningbox}

\subsection{Resumen de la sección}

El preentrenamiento en dos fases incorpora el currículo en la distribución de muestreo: primero se expande la cobertura y, posteriormente, se aumenta la densidad de señales de alta calidad. Una explicación rigurosa debe distinguir entre distribución natural, mezcla aleatoria óptima y mezcla ordenada, y especificar las métricas de mejora relativa. A continuación, el curso avanza desde los datos generales hacia los más sensibles: los datos de razonamiento (*reasoning data*).
\section{Carga anticipada del razonamiento: ¿debe el razonamiento sentar las bases durante la fase de preentrenamiento?}

Las cadenas de procesamiento tradicionales suelen describir el preentrenamiento como el aprendizaje de conocimientos generales, el ajuste fino supervisado (SFT) como la imitación de formatos y el aprendizaje por refuerzo (RL) como el verdadero fortalecimiento del razonamiento. La carga anticipada del razonamiento cuestiona esta división: si un modelo base nunca ha estado expuesto a trazas de razonamiento durante el preentrenamiento, ¿no estará obligado luego a realizar reparaciones sobre una representación débil? Lo que se estudia aquí no es «si hacer post-entrenamiento o no», sino cuántos datos de razonamiento deberían adelantarse dentro de un presupuesto fijo.

\subsection{El problema de la base débil: «razonamiento añadido demasiado tarde»}

La clase utiliza una metáfora de los cimientos para describir el flujo actual: el preentrenamiento adquiere hechos, el SFT aprende el formato superficial de las respuestas de razonamiento y el RL/RLVR es lo que finalmente fomenta la verdadera exploración. Si el razonamiento se añade únicamente como una habilidad posterior (post-hoc), la representación base podría carecer de la estructura temprana necesaria para soportar inferencias en cadena largas. Este punto de vista requiere apoyo mediante experimentos controlados posteriores y no puede sostenerse solo por medio de la metáfora.

\lecturefigure{slide-014.jpg}{El flujo tradicional coloca el razonamiento en etapas tardías, formando la hipótesis de una base débil}{Grabación oficial de Stanford, 00:12:39--00:13:29}

\begin{knowledgebox}{Concepto previo: SFT y RLVR}
El ajuste fino supervisado (SFT) utiliza forzamiento del profesor con entradas--salidas dadas o trazas de razonamiento, optimizando principalmente la probabilidad; el aprendizaje por refuerzo con recompensas verificables (RLVR) otorga recompensas a los rollout basándose en respuestas, pruebas o reglas que pueden ser evaluadas automáticamente. Ambos pertenecen al post-entrenamiento, pero difieren en las señales de optimización y los requisitos de datos.
\end{knowledgebox}

\subsection{Ejes experimentales: cuándo, cuántos, qué tan amplios y qué tan buenos}

Los experimentos sobre carga anticipada asignan el mismo tipo de datos de razonamiento a las fases de preentrenamiento, SFT y etapas posteriores, construyendo condiciones a lo largo de tres ejes: diversidad, calidad y cantidad. La base sin razonamiento no considera datos de inferencia durante el preentrenamiento, mientras que la base con razonamiento tiene acceso anticipado a una parte de ellos. Al interpretar las figuras, primero verifique si los presupuestos de datos coinciden; luego compruebe si el SFT/RL es idéntico en ambos casos; de lo contrario, «añadirlo antes» podría significar simplemente «añadir más en total».

\lecturefigure{slide-015.jpg}{Asignación de etapas de la carga anticipada del razonamiento y los tres ejes de datos}{Grabación oficial de Stanford, 00:13:29--00:14:54}

Si el presupuesto total de tokens de razonamiento es $R$, donde $R_{\mathrm{pre}}$ se asigna al preentrenamiento y $R_{\mathrm{sft}}$ al SFT, entonces
\[
R_{\mathrm{pre}}+R_{\mathrm{sft}}=R.
\]
La carga anticipada compara diferentes puntos de asignación, por ejemplo $(0,R)$ versus $(\alpha R,(1-\alpha)R)$. Si el presupuesto total o el número de rondas de entrenamiento posteriores difieren, deben marcarse como experimentos ajustados a cómputo o ajustados a datos; no pueden tratarse como la misma pregunta.

\subsection{Por qué no se deben mezclar las lecturas del modelo base y los benchmarks post-entrenados}

Un modelo base que solo ha completado el preentrenamiento y un modelo que ha completado SFT/RL suelen utilizar conjuntos de prompts, métodos de decodificación y suites de evaluación diferentes. La figura separa por etapas los conjuntos de evaluación de razonamiento general, matemáticas, ciencias y código, con el fin de advertir al lector: los promedios en páginas distintas no provienen necesariamente del mismo conjunto de preguntas; no se deben conectar directamente un 16\%, un 9.3\% y un 19\% para formar una curva monotónica.

\lecturefigure{slide-016.jpg}{Conjuntos de evaluación utilizados por el modelo base y el modelo post-entrenado}{Grabación oficial de Stanford, 00:14:54--00:15:48}

\begin{warningbox}{El error más común al interpretar las figuras}
Los benchmarks con el mismo nombre también pueden ser incomparables directamente debido a diferencias en few-shot, cadena de pensamiento (chain-of-thought), temperatura de muestreo, pass@k, extractor de respuestas o checkpoint. El material didáctico conserva las conclusiones relativas del curso, pero no interpreta los promedios entre etapas como un indicador continuo.
\end{warningbox}

\section{Cinco lecciones prácticas de cargar por adelantado}

Las siguientes cinco imágenes no demuestran repetidamente lo mismo, sino que responden secuencialmente a cinco objeciones: ¿Es efectivo inmediatamente después del preentrenamiento? ¿El SFT elimina los beneficios? ¿Sobrefitentan los datos de alta calidad? ¿Puede compensarse con más cómputo de SFT o más datos tardíos? ¿Cuánto queda después de completar RLVR? Cada gráfico debe verificar primero la comparación, luego leer la altura de las barras.

\subsection{Lección 1: Beneficios inmediatos tras el preentrenamiento}

La primera lección revisa el punto temporal más directo: antes de entrar en SFT o RL, ¿añadir datos de razonamiento altera ya al modelo base? Esta verificación elimina todos los factores posteriores al entrenamiento y observa si la asignación de datos durante el preentrenamiento es suficiente para generar una diferencia medible. Se comparan las bases con y sin razonamiento tras completar el preentrenamiento; la mejora relativa promedio en Ampere es de aproximadamente 16\%. Esto indica que los datos de razonamiento no solo enseñan al modelo a producir un formato específico, sino que pueden cambiar la representabilidad de la base para tareas matemáticas, científicas y de razonamiento general. Aún no prueba que el modelo haya generado internamente un proceso de razonamiento fielmente interpretable.

\lecturefigure{slide-017.jpg}{Beneficios inmediatos de añadir datos de razonamiento en la fase de preentrenamiento}{Grabación oficial de Stanford, 00:15:48--00:16:36}

\begin{knowledgebox}{Lectura de la figura: orden de comparación}
Verificar primero que ambos grupos de modelos son consistentes en arquitectura, presupuesto de tokens y mezcla de preentrenamiento estándar; luego observar las barras azul y verde en cada benchmark; finalmente revisar la mejora relativa promedio. Si un punto de partida es muy bajo, pequeños cambios absolutos pueden generar porcentajes relativos grandes, por lo que se debe observar tanto la puntuación cruda como la ganancia relativa.
\end{knowledgebox}

\subsection{Lección 2: El SFT no elimina las ventajas}

Que haya beneficios inmediatos no es suficiente, porque los modelos de producción generalmente continúan recibiendo SFT; por ello, la segunda lección aplica el mismo entrenamiento posterior a ambas fundaciones para verificar si la ventaja es solo un fenómeno efímero. Tras someter a dos modelos base idénticos al mismo SFT, la base con razonamiento mantiene una ventaja relativa promedio de aproximadamente 9.3\%. La conclusión pedagógica aquí es "la fundación importa": el entrenamiento posterior no reinicia completamente el modelo, sino que sigue optimizando sobre su representación base. Las ganancias tempranas pueden disminuirse, mantenerse o amplificarse en tareas específicas; se debe verificar por tarea.

\lecturefigure{slide-018.jpg}{Las ganancias iniciales de razonamiento persisten tras el SFT}{Grabación oficial de Stanford, 00:16:36--00:17:14}

\teachervoice{De 00:16:36 a 00:17:14, la oradora usa repetidamente la expresión "washed away". Le interesa no solo la puntuación base, sino si esa misma ventaja sigue siendo observable tras el entrenamiento posterior; este es el criterio clave para determinar si cargar por adelantado vale la pena desde una perspectiva de ingeniería.}

\subsection{Lección 3: Los datos de alta calidad pueden tener un efecto latente}

La tercera lección aborda un problema más sutil: si los datos de alta calidad no elevan inmediatamente las puntuaciones en la evaluación base, ¿significa que no tienen valor? El curso clasifica los datos de razonamiento en SHQ, LDQ y LMQ. SHQ es pequeño pero de alta calidad; LDQ es grande y diverso pero de menor calidad; LMQ combina ambos. Al finalizar el preentrenamiento, LMQ y LDQ pueden ser casi idénticos; tras el SFT, LMQ muestra una mejora adicional de aproximadamente 4.25\%. La señal de alta calidad no se refleja inmediatamente en los benchmarks base, sino que puede ofrecer mejor plasticidad para optimizaciones futuras.

\lecturefigure{slide-019.jpg}{El beneficio potencial de los datos de preentrenamiento de alta calidad se activa tras el SFT}{Grabación oficial de Stanford, 00:17:14--00:18:48}

\begin{knowledgebox}{Glosario: SHQ, LDQ, LMQ}
\begin{tabularx}{\textwidth}{L{0.16\textwidth}L{0.23\textwidth}X}
\toprule
Abreviatura & Atributo de los datos & Función pedagógica \\
\midrule
SHQ & Pequeño, Alta Calidad & Verificar si una cantidad pequeña de datos filtrados con precisión proporciona estructura adicional \\
LDQ & Grande, Diverso, Baja Calidad & Verificar la base de escala y diversidad \\
LMQ & Combinación de SHQ + LDQ & Verificar si el componente de alta calidad muestra un efecto latente tras el entrenamiento posterior \\
\bottomrule
\end{tabularx}
\end{knowledgebox}

\subsection{Lección 4: Duplicar cómputo tardío o desplazar datos no logra alcanzarlo}

En la figura izquierda se le dan más épocas de SFT a la base sin razonamiento; en la derecha se fija el presupuesto total de datos de razonamiento y se compara "todo en SFT" frente a "preentrenamiento + asignación en SFT". En ambos casos, el modelo cargado por adelantado mantiene su ventaja. Esto refuta respectivamente las explicaciones de que "solo falta más entrenamiento posterior" o que "basta con usar los mismos datos más tarde".

\lecturefigure{slide-020.jpg}{Efecto fuerte de la fundación bajo condiciones de cómputo y datos emparejados}{Grabación oficial de Stanford, 00:18:48--00:20:54}

\begin{importantbox}{Por qué realizar ambos tipos de emparejamiento simultáneamente}
El emparejamiento de cómputo controla la cantidad de entrenamiento computacional; el emparejamiento de datos controla el presupuesto de muestras; ambos resuelven factores de confusión distintos. Si la condición temprana usa más datos y consume más FLOPs, no se puede atribuir el beneficio únicamente al momento. Un reporte riguroso debe listar al menos los tokens, épocas, pasos del optimizador, longitud de secuencia y el costo adicional de rollout/SFT.
\end{importantbox}

\subsection{Lección 5: Persiste una ventaja duradera tras RLVR}

Que siga liderando después del SFT no garantiza que RLVR no vuelva a nivelar ambas trayectorias; la quinta lección verifica entonces el pipeline final entregable, en lugar de detenerse en un checkpoint intermedio. La ventaja relativa promedio de Ampere frente a Volta es de aproximadamente 19\%, con diferencias relativas aún mayores en tareas matemáticas difíciles como AIME. El gráfico apoya la idea de que "el razonamiento temprano puede generar una ventaja duradera y acumulativa", pero no indica que todas las tareas crezcan al mismo ritmo, ni que el entrenamiento posterior deje de ser importante; es precisamente el entrenamiento posterior idéntico lo que expone las diferencias entre fundaciones.

\lecturefigure{slide-021.jpg}{Cargar por adelantado mantiene la ventaja tras el entrenamiento completo posterior}{Grabación oficial de Stanford, 00:20:54--00:21:48}

\begin{warningbox}{Una mejora relativa grande no implica cercanía a la puntuación máxima}
Cuando las tareas como AIME tienen una línea base baja, varias puntos de mejora absoluta pueden corresponder a una ganancia relativa muy grande. Las decisiones de ingeniería deben reportar simultáneamente la puntuación absoluta, el promedio, el número de ensayos y la varianza; no basta con recortar el porcentaje relativo más alto para afirmar una mejora general de capacidades.
\end{warningbox}

\subsection{¿Qué detalles de ingeniería fuera de los artículos completó la sesión de preguntas y respuestas?}

La sesión de preguntas intermedias aclaró que el uso de datos de razonamiento se basa en conjuntos de datos existentes en la comunidad y definiciones de trabajo comunes, donde las muestras típicas son pregunta, rastro de razonamiento largo y respuesta final, con un sesgo experimental hacia STEM. La calidad puede operacionalizarse mediante intensidad de filtrado, complejidad, longitud o etiqueta de dificultad; la optimización de la receta es un proceso híbrido de estimación automática y ensamblaje manual. La oradora también especificó que invertir el orden de las fases resultó peor en sus estudios de abstracción.

\teachervoice{De 00:22:00 a 00:27:58, la oradora no empaquetó "razonamiento" como un concepto unificado transversal. Reconoció que los datos usados se inclinan hacia matemáticas, código y STEM, y el blend final tampoco es totalmente automático. Estas limitaciones determinan los límites de extrapolación de los resultados.}

\begin{knowledgebox}{Práctica profesional: cómo registrar una receta de datos reproducible}
Se debe guardar al menos la versión de la fuente de datos, las claves de deduplicación, filtros y umbrales, rúbrica de calidad, mapeo de dominios, número de tokens únicos, épocas objetivo, pesos por fase, límite de fases, semilla aleatoria y configuración de evaluación proxy. Publicar solo el porcentaje final no permite distinguir diferencias causadas por estimaciones de calidad, implementación de muestreo o orden.
\end{knowledgebox}

\subsection{Resumen de la sección}

Las cinco lecciones establecen la cadena de evidencia para cargar por adelantado: los beneficios aparecen inmediatamente, no desaparecen tras el SFT, las señales de alta calidad pueden manifestarse con retraso, más cómputo o datos tardíos no logran igualar completamente y la ventaja persiste incluso después del RLVR completo. La siguiente parte profundiza aún más: además de proporcionar datos de razonamiento antes, ¿podemos hacer que el modelo explore activamente sus propios pensamientos durante el preentrenamiento?
\section{Motivación de RLP: del texto observacional al aprendizaje mediante pensamiento}

RLP ya no solo ajusta el orden de los datos, sino que introduce una acción pensable (thought action) muestreable antes de cada token seleccionado. La idea central es: el modelo genera primero varias trazas de razonamiento interno y luego predice el siguiente token real; si cierta thought aumenta la probabilidad de ese token, recibe una recompensa positiva. De este modo, el texto de preentrenamiento estándar proporciona simultáneamente el token objetivo y una referencia de recompensa auto-supervisada.

\subsection{La historia de dos aprendices}

Leo primero construye un puente rudimentario pero funcional para un vehículo pequeño, luego lo mejora a partir de la retroalimentación; Bolt lee todo el material sobre puentes y crea modelos hermosos que no cumplen las restricciones. La historia destaca la diferencia entre aprender haciendo (learning by doing) y aprendiendo observando (learning by observing). La predicción del siguiente token es hábil imitando patrones estadísticos a partir de textos ya escritos, mientras que RLP intenta otorgar al modelo una oportunidad de exploración para «proponer primero una explicación y luego verificarla con el token real».

\lecturefigure{slide-023.jpg}{Diferencia entre dos aprendices: ensayo y error práctico versus observación de ejemplos}{Grabación oficial de Stanford, 00:28:14--00:29:37}

\begin{warningbox}{Thought no es una explicación automáticamente confiable}
Las cadenas de pensamiento (chain-of-thought) generadas por el modelo pueden contener racionalizaciones a posteriori, hechos erróneos o atajos incomprensibles. RLP solo verifica si la thought aumenta la probabilidad de predicción del siguiente token observado; no valida si coincide con el razonamiento humano, si es causalmente fiel o si se mantiene fuera de la distribución.
\end{warningbox}

\subsection{El problema del preentrenamiento estándar no es la falta de gradientes, sino la ausencia de acciones de exploración}

El NTP estándar ya proporciona gradientes densos, pero el espacio de acciones (action space) solo permite ajustes directos a la distribución del siguiente token; el razonamiento complejo suele aparecer explícitamente recién en SFT, RLHF o RLVR. La formulación del problema de RLP es: ¿es posible anticipar la exploración manteniendo al mismo tiempo la escalabilidad del flujo documental estándar? Esta motivación también responde al muro de datos: si los nuevos tokens se vuelven cada vez más costosos, hay que aumentar la señal de aprendizaje generada por cada token.

\lecturefigure{slide-024.jpg}{El preentrenamiento estándar retrasa la exploración explícita hasta el post-entrenamiento}{Grabación oficial de Stanford, 00:29:37--00:31:23}

El objetivo del NTP estándar es
\[
\mathcal{L}_{\mathrm{NTP}}(\theta)
=-\sum_{t=1}^{L}\log p_\theta(x_t\mid x_{<t}).
\]
Donde, $x_{<t}$ representa el contexto, $x_t$ es el siguiente token real en el corpus y $p_\theta$ es la distribución del modelo. Este objetivo aprende mucha estructura, pero no tiene una variable separada que represente «qué thought se exploró antes de generar la salida».

\subsection{Diferencia entre Vanilla y RLP en probabilidades condicionales}

En el ejemplo de la fotosíntesis, el modelo vanilla predice directamente `sunlight` a partir del contexto; el modelo RLP genera primero una thought que dice «este proceso depende de energía solar», luego condicionalmente al contexto y a esa thought predice el mismo token. La clave no es generar más texto, sino tratar la thought como una acción optimizable y comparar su ganancia relativa respecto a un baseline que solo usa el contexto.

\lecturefigure{slide-025.jpg}{Predicción vanilla del siguiente token versus predicción RLP condicionada al pensamiento}{Grabación oficial de Stanford, 00:31:23--00:32:34}

\[
p_\theta(x_t\mid x_{<t},c_t)
\quad\text{versus}\quad
p_\phi(x_t\mid x_{<t}),
\]
donde, $c_t$ es la thought muestreada antes de la posición $t$, $\theta$ son los parámetros del policy de thoughts o predictor razonado, y $\phi$ son los parámetros del baseline sin pensamiento. Ambos ven el mismo contexto y el mismo token objetivo real; la diferencia radica únicamente en si se incluye $c_t$.

\begin{importantbox}{Ciclo auto-supervisado de RLP}
El token real $x_t$ del corpus cumple simultáneamente dos funciones: es el objetivo para NTP y sirve como referencia para evaluar si la thought tiene utilidad predictiva. Por ello, RLP no requiere un verificador de respuestas manual, pero sí depende del siguiente token real; no recompensa de forma insupervisada cualquier texto que «parezca razonamiento».
\end{importantbox}

\subsection{Resumen de la sección}

RLP convierte la thought en una acción exploratoria antes de la predicción y verifica si dicha acción aporta información mediante el siguiente token real. La diferencia con NTP no es si existen señales de entrenamiento, sino si hay un camino de exploración explícito y recompensas relativas sobre la calidad del camino. A continuación se desglosa paso a paso el ciclo de entrenamiento completo.
\section{Mecanismo RLP: política de pensamiento, línea base contrafactual y ganancia de información}

El pipeline completo de RLP consta de cuatro partes interdependientes: la política de pensamiento genera el rollout, la línea base sin pensar proporciona predicciones contrafactuales, la recompensa de ganancia de información mide el beneficio del pensamiento y el EMA (Promedio Móvil Exponencial) permite que la línea base siga de manera estable a la política. Mirar solo la fórmula final hace que se pierdan las condiciones de ingeniería como el prompt, el muestreo, los gradientes de parada y la cronología de las actualizaciones.

\subsection{Primero observe dos rutas en el diagrama de flujo global}

Una entrada del contexto entra en la línea base sin pensar, obteniendo $p_\phi(x_t\mid x_{<t})$; la otra ruta entra en la política de pensamiento, primero se muestrea $c_t$, luego se obtiene $p_\theta(x_t\mid x_{<t},c_t)$. La diferencia entre dos probabilidades log (log probability) forma la recompensa, y esta recompensa actualiza la política de los tokens de pensamiento, mientras que el NTP estándar sigue entrenando al modelo para predecir el token real.

\lecturefigure{slide-026.jpg}{Visión general del RLP: dos rutas de predicción, recompensa de ganancia de información y actualización de políticas}{Grabación oficial de Stanford, 00:32:34--00:33:00}

\begin{knowledgebox}{Lectura de la figura: no interprete las dos redes como sistemas permanentemente independientes}
En el diagrama, $\theta$ y $\phi$ representan dos estados de parámetros: la política es la red actual y la línea base es su copia con retraso en EMA. Comparten la misma arquitectura; la línea base no es un verificador externo entrenado por separado. La réplica con retraso proporciona una referencia contrafactual que cambia lentamente.
\end{knowledgebox}

\subsection{El contrato de prompt define el espacio de exploración}

El prompt mostrado en clase requiere que el modelo se centre en el siguiente paso de razonamiento, sin saltar directamente a la respuesta final en caja ni repetir la pregunta u outputear metadatos. Esto indica que "hacer pensar al modelo" no es neutral: el prompt define la longitud del pensamiento, el estilo y el espacio de exploración. Si el prompt cambia, también cambian la distribución de recompensas y el cómputo del rollout.

\lecturefigure{slide-027.jpg}{Contrato de prompt para generar thought}{Grabación oficial de Stanford, 00:33:00--00:33:19}

\begin{lstlisting}[caption={Versión simplificada de enseñanza del prompt RLP thought}]
Prompt:
  Razone sobre el siguiente paso útil para predecir el próximo token.
  No salte a una respuesta final en caja.
  No repita la pregunta, los metadatos del contexto ni los comentarios.
  Produzca un pensamiento conciso que pueda ayudar a la predicción del siguiente token.
\end{lstlisting}

\begin{warningbox}{El prompt también es un hiperparámetro del algoritmo}
La longitud del pensamiento, la temperatura de muestreo, el número de rollouts, el token de parada y la redacción del prompt alteran el costo de exploración y las señales aprendibles. Omitir estos detalles e informar solo que se "usa chain-of-thought" hará que los resultados de replicación pierdan comparabilidad.
\end{warningbox}

\subsection{Política de pensamiento: primero generar acción, luego predecir token}

Dado un prompt y un contexto, la política muestrea múltiples rollouts de pensamiento. Cada pensamiento concatenado con el contexto predice el siguiente token real; por tanto, un solo rollout genera simultáneamente la probabilidad de acción y la probabilidad del token. La optimización relativa al grupo puede comparar múltiples pensamientos en la misma posición entre sí, reduciendo la necesidad de entrenar un modelo de valor separado.

\lecturefigure{slide-028.jpg}{Generación de rollout de razonamiento y predicción del siguiente token mediante la política de pensamiento}{Grabación oficial de Stanford, 00:33:19--00:34:02}

Para el $i$-ésimo rollout, denotemos el pensamiento como $c_t^{(i)}$; entonces su puntuación razonada es
\[
S_{\mathrm{pred}}^{(i)}
=\log p_\theta(x_t\mid x_{<t},c_t^{(i)}).
\]
Donde, $i=1,\ldots,G$, $G$ es el número de pensamientos muestreados en cada posición. Una $S_{\mathrm{pred}}$ alta no significa que el texto del pensamiento sea correcto por sí mismo, sino que indica que ayuda al modelo a asignar mayor probabilidad al token real.

\subsection{Línea base sin pensar: se requiere un contrafactual con el mismo contexto}

Si solo maximizamos la puntuación razonada, el modelo podría obtener altas puntuaciones aprovechando tokens que de todas formas son fáciles. La línea base utiliza el mismo contexto pero no considera el pensamiento, estimando lo fácil que ya era ese token "sin exploración". Por tanto, la recompensa se centra en el incremento aportado por el pensamiento y no en la verosimilitud absoluta del token.

\lecturefigure{slide-029.jpg}{La línea base sin thinking proporciona un contrafactual basado solo en contexto}{Grabación oficial de Stanford, 00:34:02--00:34:25}

\begin{importantbox}{Función del contrafactual}
Si un mismo $x_t$ era prácticamente seguro por sí mismo, entonces $p_\phi$ ya es alto y un pensamiento común difícilmente obtendría una gran recompensa; si algún pensamiento elimina realmente la ambigüedad, el aumento relativo de $p_\theta$ frente a $p_\phi$ será significativo. La línea base separa la "simplidad del problema" de la "eficacia del razonamiento".
\end{importantbox}

\subsection{Ganancia de información: relación log-verosimilitud}

La recompensa central del RLP es la diferencia en la probabilidad log entre el predictor razonado y la línea base sin pensar para el token real. Un valor positivo indica que el pensamiento mejora la predicción, cero significa que no contribuye y un valor negativo indica que el pensamiento engaña al modelo. El artículo trata la recompensa como una constante al calcular la actualización de la política, evitando propagar los gradientes a través de ambos evaluadores para evitar un acoplamiento descontrolado del objetivo.

\lecturefigure{slide-030.jpg}{Diagrama completo del cálculo de la recompensa de ganancia de información}{Grabación oficial de Stanford, 00:34:25--00:35:34}

\[
r_t^{(i)}
=\log p_\theta(x_t\mid x_{<t},c_t^{(i)})
-\log p_\phi(x_t\mid x_{<t}).
\]
Donde, $r_t^{(i)}$ es la recompensa del $i$-ésimo pensamiento en la posición $t$; $p_\theta$ es la probabilidad de predicción tras añadir el pensamiento; $p_\phi$ es la línea base basada solo en contexto; $x_t$ es siempre el token real del corpus. Esta magnitud es similar a una relación log-verosimilitud punto por punto, no a la esperanza completa del conjunto de datos de la información mutua de Shannon.

\teachervoice{00:34:25--00:35:34, el orador enfatiza: la recompensa solo es positiva cuando el pensamiento "contribuye significativamente"; los pensamientos basura pueden obtener cero o valores negativos. Esta aclaración verbal se acerca más a la esencia del método que decir simplemente "el modelo genera CoT".}

\begin{warningbox}{La ganancia de información no equivale a la verdad}
Un pensamiento erróneo pero estadísticamente relacionado con el corpus de entrenamiento también podría aumentar la probabilidad de algún token; un pensamiento correcto pero redundante o con redacción inadecuada podría tener una recompensa muy baja. El RLP optimiza la utilidad predictiva, sin optimizar directamente la veracidad (truthfulness), fidelidad o seguridad.
\end{warningbox}

\subsection{Recompensa densa no binaria y ventaja relativa al grupo}

La ganancia de información define una puntuación para un solo pensamiento, pero la optimización de la política también necesita convertir múltiples rollouts en crédito relativo; esta sección transita desde la recompensa hacia la ventaja, explicando por qué el RLP no necesita entrenar un modelo de valor adicional para cada token. Diferente a las recompensas escasas que dan 0/1 solo en la respuesta final, la recompensa del RLP es un número real y puede calcularse en cualquier posición seleccionada del documento. Para los $G$ rollouts en una misma posición, se puede construir una ventaja usando el promedio y la desviación estándar dentro del grupo, otorgando actualizaciones positivas a aquellos rollouts que "ayudan más que otros pensamientos del mismo grupo".

\lecturefigure{slide-031.jpg}{El RLP utiliza recompensa token-level densa y no binaria}{Grabación oficial de Stanford, 00:35:34--00:35:45}

\[
A_t^{(i)}
=\frac{r_t^{(i)}-\bar r_t}{\operatorname{std}(r_t)+\epsilon},
\qquad
\bar r_t=\frac{1}{G}\sum_{i=1}^{G}r_t^{(i)}.
\]
Donde, $A_t^{(i)}$ es la ventaja relativa al grupo, $G$ es el número de rollouts y $\epsilon$ evita división por cero cuando la varianza es muy pequeña. El RLP real utiliza una actualización proxy recortada (clipped surrogate) para los tokens de pensamiento y entrelaza esto con el entrenamiento estándar de verosimilitud.

\subsection{Línea base EMA: debe seguir sin perseguir al unísono}

Si la línea base está completamente congelada, pierde gradualmente su valor de referencia; si por el contrario se sincroniza completamente con la política en cada paso, podría hacer que el objetivo de recompensa se desvíe rápidamente. El Promedio Móvil Exponencial (EMA) ofrece un compromiso con una actualización lenta: la línea base mantiene lo "suficientemente nuevo", pero intencionalmente queda atrás, mitigando así el hacking de recompensas y el arranque autoinestable.

\lecturefigure{slide-032.jpg}{Ubicación de la actualización de la línea base no-think con retraso en EMA}{Grabación oficial de Stanford, 00:35:45--00:36:31}

\[
\phi\leftarrow \tau\phi+(1-\tau)\theta,
\qquad 0<\tau<1.
\]
Donde, $\theta$ son los parámetros de la política de pensamiento actual y $\phi$ son los parámetros de la línea base; $\tau$ es el decaimiento del EMA. Cuanto más cercano esté $\tau$ a 1, más estable pero antigua será la línea base; cuanto menor sea, más rápido seguirá, aunque el objetivo contrafactual será más propenso a oscilar.

\begin{lstlisting}[caption={Pseudocódigo simplificado para un paso de entrenamiento del RLP}]
document = sample_general_pretraining_document()
t = random_token_position(document)
context, target = document[:t], document[t]

thoughts = policy.sample_thoughts(context, group_size=G)
baseline_logp = no_think_ema.log_prob(target, context)
rewards = []
for thought in thoughts:
    reasoned_logp = policy.log_prob(target, context, thought)
    rewards.append(stop_gradient(reasoned_logp - baseline_logp))

advantages = group_normalize(rewards)
update_thought_policy_clipped(thoughts, advantages)
update_standard_next_token_loss(document)
ema_update(no_think_ema, policy, decay=tau)
\end{lstlisting}

\subsection{Resumen de la sección}

El ciclo cerrado del RLP es: la política muestrea pensamientos, la línea base ofrece un contrafactual sin pensar, la relación log-verosimilitud forma una recompensa densa, la ventaja relativa al grupo actualiza los tokens de pensamiento, el EMA estabiliza la línea base y, simultáneamente, el NTP sigue aprendiendo del documento. Si falta cualquier componente, el método podría degenerarse en una imitación común de CoT o en un autorecompensado inestable.

\section{Experimentos de RLP: cuatro dimensiones para una comparación equitativa}

RLP requiere rollout; por ello, el costo computacional por token de entrenamiento es mayor que en NTP. La evaluación equitativa no puede limitarse a comparar la cantidad de tokens; también debe compararse el número de FLOPs, el checkpoint inicial, los SFT/RL posteriores, la arquitectura del modelo y los benchmarks. La clase presenta secuencialmente Qwen3-1.7B, una línea base equivalente en cómputo (compute-equivalent baseline) y Nemotron-Nano-12B-V2, formando así una cadena de evidencia que va desde modelos pequeños hasta arquitecturas híbridas.

\subsection{Qwen3-1.7B: las condiciones del experimento preceden al gráfico de barras}

El primer conjunto de experimentos parte del final checkpoint de Qwen3-1.7B-Base para realizar RLP sobre un corpus de preentrenamiento general, sin exigir un conjunto de datos específico de razonamiento. La línea base incluye la versión base original, el preentrenamiento continuado (CPT) con el mismo número de tokens y las versiones que recibieron los mismos procesos de SFT/RLVR. Solo tras confirmar estas condiciones se puede explicar si la "fundación" mantiene sus capacidades tras el post-entrenamiento.

\lecturefigure{slide-033.jpg}{Condiciones de entrenamiento, línea base y post-entrenamiento en los experimentos de RLP con Qwen3-1.7B}{Grabación oficial de Stanford, 00:36:31--00:37:34}

\begin{knowledgebox}{Términos clave: CPT y RLP}
El preentrenamiento continuado (CPT) sigue utilizando la pérdida estándar de NTP, limitándose a procesar más tokens sobre el checkpoint base; en cambio, RLP destina parte del cómputo al rollout de pensamientos y a actualizaciones basadas en ganancia de información. Se requiere tanto CPT igualado por tokens como CPT igualado por FLOPs para un control adecuado.
\end{knowledgebox}

\subsection{Resultados con coincidencia de tokens y persistencia tras el post-entrenamiento}

Bajo el mismo número adicional de tokens, RLP mejora el rendimiento promedio respecto tanto a la base como al CPT; dicha ventaja perdura incluso después de completar el mismo post-entrenamiento. Dado que los benchmarks en los lados izquierdo y derecho del gráfico corresponden a etapas diferentes, no se deben restar directamente las alturas de las barras; el énfasis pedagógico radica en la consistencia de la dirección: la fundación de razonamiento establecida por RLP no ha sido eliminada por el alineamiento.

\lecturefigure{slide-034.jpg}{Resultados de token-matched RLP y post-entrenamiento en Qwen3-1.7B}{Grabación oficial de Stanford, 00:37:34--00:38:26}

\begin{knowledgebox}{Lectura de la figura: primero diferenciar lados, luego colores}
El lado izquierdo compara modelos de base-stage; el derecho compara modelos que recibieron los mismos procesos de SFT/RLVR. El color verde indica la ruta con RLP y el azul su línea base correspondiente. Es necesario observar Math, Science, Science Post y Overall simultáneamente, en lugar de seleccionar solo la barra más alta. La figura respalda un aumento relativo promedio dentro del suite de benchmarks, pero no demuestra que todas las áreas de competencia crezcan proporcionalmente.
\end{knowledgebox}

\subsection{Coincidencia de FLOPs: el rollout no es gratuito}

RLP genera múltiples pensamientos por cada posición objetivo y recalcula la probabilidad condicionada al pensamiento; por ello, el costo de 1B tokens RLP difiere del de 1B tokens CPT. La clase asocia los 170M tokens RLP a aproximadamente 6B tokens de CPT igualado en FLOPs, lo que equivale a que el CPT procese unas 35 veces más datos. Aún así, RLP conserva una ventaja relativa promedio de aproximadamente 14\% en la configuración reportada.

\lecturefigure{slide-035.jpg}{Contraste con CPT igualado en FLOPs en Qwen3-1.7B}{Grabación oficial de Stanford, 00:38:26--00:40:19}

El cálculo aproximado se puede expresar como
\[
C_{\mathrm{RLP}}
\approx C_{\mathrm{NTP}}+G\,C_{\mathrm{rollout}}+G\,C_{\mathrm{score}},
\]
donde $G$ es el número de rollouts de pensamientos, $C_{\mathrm{rollout}}$ es el costo de generar los pensamientos y $C_{\mathrm{score}}$ es el costo de puntuar los tokens objetivo. La contabilidad real de FLOPs también depende del reutilización de KV, la longitud de los pensamientos, el batching, el empaquetado de secuencias y la implementación de kernels.

\begin{warningbox}{La eficiencia en tokens no es lo mismo que la eficiencia en tiempo de reloj}
Obtener una puntuación más alta con menos tokens de datos mediante RLP no implica que el entrenamiento sea más rápido o más barato. La generación secuencial del rollout, una utilización menor del hardware, las activaciones adicionales y la comunicación pueden incrementar tanto el tiempo como el consumo energético. Las decisiones de despliegue deben reportar simultáneamente tokens, FLOPs, horas GPU y throughput.
\end{warningbox}

\subsection{Nemotron-Nano-12B-V2: modelos más grandes y arquitectura híbrida}

El segundo conjunto de experimentos parte de un checkpoint intermedio de aproximadamente 19.8T tokens, añadiendo solo 250M tokens RLP para compararse con el final base entrenado hasta unos 20T tokens. Nemotron-Nano-12B-V2 es una arquitectura híbrida Mamba-Transformer; por ello, este experimento prueba simultáneamente la escala del modelo, la etapa inicial y la migración de arquitectura. El control clave es que la ruta RLP consume aproximadamente 200B tokens menos en total.

\lecturefigure{slide-036.jpg}{Configuración de RLP con datos limitados sobre checkpoint intermedio de Nemotron}{Grabación oficial de Stanford, 00:40:19--00:41:14}

\teachervoice{00:40:19--00:42:20, el orador enfatiza que la ruta RLP parte del checkpoint de 19.8T y utiliza solo 250M tokens, mientras que su objeto de comparación es el final base de 20T. Este diseño busca demostrar que los beneficios no provienen de "ver más datos", sino de cambiar los objetivos de aprendizaje con menos tokens en total.}

La página de resultados muestra un aumento relativo promedio de aproximadamente 35\% en la etapa base, siendo especialmente notable el incremento en el dominio science; tras el mismo post-entrenamiento se mantiene una ventaja absoluta. La figura respalda la portabilidad del RLP a través de modelos Transformer y híbridos, pero dado que los modelos y datos cubren solo casos limitados, no se puede afirmar que todas las arquitecturas presenten la misma pendiente de escalado.

\lecturefigure{slide-037.jpg}{Resultados de base y post-entrenamiento del RLP en un modelo híbrido Nemotron más grande}{Grabación oficial de Stanford, 00:41:14--00:42:20}

\begin{knowledgebox}{Lectura de la figura: no confundir el relativo 35\% con el absoluto 3\%}
El lado izquierdo aborda el aumento relativo promedio en la etapa base; el derecho compara las barras post-entrenadas usualmente mediante puntuaciones absolutas. Ambas métricas responden preguntas distintas: la primera examina la proporción de mejora de la fundación, mientras que la segunda indica cuánto sigue liderando el modelo final utilizable. Al reportar los resultados es obligatorio mantener el denominador y el suite de benchmarks.
\end{knowledgebox}

\subsection{Resumen de la sección}

La credibilidad de los experimentos de RLP surge de múltiples líneas base: la coincidencia de tokens verifica la cantidad de datos, la coincidencia de FLOPs valida el costo del rollout, la coincidencia de post-entrenamiento comprueba si los beneficios son duraderos y la comparación de checkpoints intermedios examina si es posible superar al final base con menos tokens en total. El próximo capítulo situará a RLP dentro del corpus literario del razonamiento temprano.

````
\section{Early Reasoning: Diferencias entre las referencias RLP, Quiet-STaR, RPT y RLPT}

Hacer que el modelo "piense antes de hablar" no fue una innovación exclusiva de RLP. Tanto Quiet-STaR como Reinforcement Pre-Training (RPT) y Reinforcement Learning on Pre-Training Data (RLPT) anticipan el razonamiento. Lo verdaderamente importante para comparar es dónde proviene la recompensa, en qué nivel de granularidad se calcula, si requiere un verificador externo y si la trazabilidad del razonamiento obtiene crédito directamente.

\subsection{De Quiet-STaR al preentrenamiento por refuerzo}

La página de las referencias coloca a Quiet-STaR, RPT y RLPT junto con RLP en una misma línea temporal. Quiet-STaR genera y evalúa el pensamiento dentro del texto ordinario; RPT convierte la corrección del siguiente token en una recompensa de aprendizaje por refuerzo verificable; RLPT utiliza datos de preentrenamiento para construir tareas de refuerzo a nivel de fragmentos más largos; mientras que RLP usa la ganancia en log-likelihood derivada del pensamiento como señal intrínseca densa para el siguiente token real.

\lecturefigure{slide-038.jpg}{Ruta de early reasoning: Quiet-STaR, RPT, RLPT y RLP}{Grabación oficial de Stanford, 00:42:20--00:43:13}

\begin{knowledgebox}{Tabla de términos: qué resuelven los métodos adyacentes}
\begin{tabularx}{\textwidth}{L{0.18\textwidth}X X}
\toprule
Método & Mecanismo central & Diferencia clave con RLP \\
\midrule
NTP & Maximización directa de la verosimilitud del siguiente token real & Sin acción de pensamiento explícita / recompensa \\
Quiet-STaR & Muestreo de pensamientos en el texto y medición de su utilidad & Enfoque temprano de auto-pensamiento, los detalles de entrenamiento difieren de RLP \\
RPT & Convertir el razonamiento del siguiente token en una tarea de RL verificable & Uso de filtrado y recompensa binaria escasa (sparse) de corrección \\
RLPT & Realizar RL a nivel de segmento sobre datos de preentrenamiento & Diferencias en la granularidad de la recompensa y el diseño del verificador \\
RLP & Recompensar el rollout con la ganancia de log-likelihood aportada por el pensamiento & Densa, a nivel de token, contrafactual sin pensar (no-think) \\
\bottomrule
\end{tabularx}
\end{knowledgebox}

\subsection{Tres ejes de comparación: fuente de recompensa, granularidad y emergencia}

NTP no tiene una recompensa separada; RPT/RLPT utilizan señales externas o verificables, generalmente más escasas; la recompensa de RLP proviene de la diferencia entre el predictor razonado interno y la línea base EMA, tomando valores continuos en la posición del token. El curso etiqueta a RLP como emergencia explícita/fuerte de razonamiento, lo cual es una generalización sobre las señales de entrenamiento por parte del autor, no una medición directa de los procesos cognitivos internos.

\lecturefigure{slide-039.jpg}{Comparación de recompensa y razonamiento entre NTP, RPT/RLPT y RLP}{Grabación oficial de Stanford, 00:43:13--00:43:56}

\begin{warningbox}{Definición precisa de Verifier-free}
RLP no requiere un modelo de determinación de respuestas adicional ni un verificador con anotaciones humanas; sin embargo, sigue utilizando el siguiente token real del corpus como supervisión y emplea una línea base EMA. Por tanto, es más preciso decir "libre de verificador externo" en lugar de "sin supervisión o sin referencia".
\end{warningbox}

\subsection{Comparación con RPT emparejada: leer junto con el mecanismo y los números}

Una tabla cualitativa solo ilustra las diferencias de diseño, pero no indica qué método es más eficaz bajo presupuestos similares; por ello, el curso sigue comparando RPT y RLP utilizando un entorno emparejado (matched setting) que sitúa la mecánica de recompensa y los resultados cuantificados en la misma página. Bajo configuraciones similares como Qwen3-1.7B-Base, OmniMath y 170 millones de tokens, el curso reporta que RLP supera a RPT con un promedio del 4\%. La explicación del mecanismo es: RPT primero utiliza un filtro externo para encontrar posiciones viables y luego otorga una recompensa binaria escasa; en cambio, RLP puede calcular la ganancia de información densa en cualquier token seleccionado y vincular el crédito a la trazabilidad del razonamiento.

\lecturefigure{slide-040.jpg}{Comparación en entorno emparejado entre RLP y RPT}{Grabación oficial de Stanford, 00:43:56--00:45:56}

\begin{knowledgebox}{Lectura de la figura: qué variables aún pueden influir en la comparación}
Aunque el modelo, los datos y el número de tokens sean similares, el prompt, el número de rollouts, la longitud del pensamiento, el optimizador, el recorte KL/clipping, el filtrado de muestras y el cómputo pueden diferir. La figura respalda que "RLP funciona mejor en los experimentos replicados por los autores", pero esto no equivale a una dominación teórica sobre todas las implementaciones de RPT.
\end{knowledgebox}

\subsection{Resumen de la sección}

RLP pertenece a la familia de early reasoning / preentrenamiento por refuerzo; su distintividad proviene del contrafactual sin pensar (no-think counterfactual) y de la recompensa basada en ganancia de información densa. Al comparar métodos adyacentes, es obligatorio examinar simultáneamente la fuente de recompensa, la granularidad, los costos de filtrado y el cómputo, en lugar de limitarse a comparar nombres de métodos o promedios individuales.
\section{Conclusiones de ingeniería, limitaciones y preguntas del curso}

El último grupo de diapositivas condensa los resultados de la investigación en principios de ingeniería: recompensar los pensamientos intermedios, usar señales densas y verificar la eficiencia de los tokens, combinando además el currículo, el razonamiento front-loading y el reforzamiento. Sin embargo, las conclusiones realmente aplicables requieren pasar por auditoría de datos, contabilidad de costos, análisis de fallos y límites de seguridad.

\subsection{Tres ablationes: por qué funciona, cuán densas deben ser las señales y cuántos datos ahorrar}

La página de Key Ablations desglosa el valor del RLP en tres puntos: Reward to Think verifica si los pensamientos obtienen crédito directo; Sparse vs Dense examina la diferencia entre recompensar solo el resultado correcto final y la ganancia de información por posición; Token Efficiency evalúa la utilización de datos relativa a una línea base NTP enorme. Las tres columnas demuestran conjuntamente que los beneficios provienen del diseño del objetivo, no simplemente de generar más segmentos de CoT.

\lecturefigure{slide-041.jpg}{Conclusiones de reward-to-think, recompensa densa y eficiencia de tokens del RLP}{Grabación oficial de Stanford, 00:45:56--00:46:27}

\begin{warningbox}{Las ablationes aún requieren una tabla de costos completa}
Si la recompensa densa utiliza más rollouts, más forwards de scorer o pensamientos más largos, también altera la señal y el cómputo. Una ablación ideal debería reportar los mismos FLOPs, el mismo tiempo wall-clock, el mismo número de tokens o varios métodos de emparejamiento, evitando confundir el cómputo adicional con el efecto real de la forma de recompensa.
\end{warningbox}

\subsection{De Pascal a Hopper: las capacidades provienen de la superposición de estrategias}

La curva reconecta los cuatro learners: el currículo aporta la primera capa de mejora, el front-loading del razonamiento eleva más el fundamento y, al agregar RLP con learning through thinking, se alcanza la mayor mejora relativa. La curva es una síntesis pedagógica a través de experimentos, no una descomposición causal estrictamente aditiva en un único entrenamiento unificado; por lo tanto, no se deben tomar los porcentajes de cada tramo como constantes independientes y superponibles.

\lecturefigure{slide-042.jpg}{De Pascal a Hopper: agregar gradualmente tres estrategias de aprendizaje}{Grabación oficial de Stanford, 00:46:27--00:47:12}

\begin{importantbox}{Juicio integral en lugar de suma mecánica}
El currículo cambia la distribución que ve el modelo, el front-loading altera la representación base y RLP modifica nuevamente el paisaje de optimización de la política de pensamiento; existen interacciones entre ellos. En una verdadera combinación se debe realizar una ablación factorial, no sumar los headline gains de cada artículo para predecir el beneficio final.
\end{importantbox}

\subsection{Seis conclusiones y tres no-conclusiones}

El curso concluye finalmente: dos fases son efectivas; el front-loading genera una ventaja duradera y compuesta; las señales de alta calidad pueden manifestarse después del SFT; RLP anticipa el razonamiento por reforzamiento; los textos sin etiquetar estándar también pueden proporcionar señales de entrenamiento tipo razonamiento; y la exploración junto con la recompensa abren un nuevo eje para el escalado del razonamiento. Cada una debe ir acompañada de límites de modelo, datos y evaluación.

\lecturefigure{slide-043.jpg}{Currículo, razonamiento y reforzamiento redefinen el preentrenamiento}{Grabación oficial de Stanford, 00:47:12--00:49:00}

En la figura \textbf{no se demuestra} nada de esto: en primer lugar, que todas las cadenas de pensamiento merecen ser retenidas; en segundo, que RLP reduciría automáticamente las alucinaciones o mejoraría el alineamiento; y en tercer lugar, que los resultados solo de lenguaje transferirían directamente a un modelo de visión-lenguaje. El orador aclara explícitamente en Q\&A que multimodal se considera no testeado, mientras que RLHF/alucinaciones son temas adyacentes pero no parte de este trabajo.

\subsection{La receta final vuelve a la colaboración sistémica}

La página de cierre muestra nuevamente los cuatro componentes, pero ahora cada columna tiene un significado concreto: Smart Data requiere etiquetas de calidad/dominio rastreables; Smart Architecture necesita soporte para entrenamiento largo y rollouts; Smart Algorithms incluye currículo, front-loading y RLP; Smart Collaboration se encarga de unificar las métricas de preentrenamiento/post-entrenamiento y el presupuesto de costos. Recuerda al lector que los artículos de algoritmos solo pueden convertirse en capacidades del modelo al ingresar a un pipeline completo.

\lecturefigure{slide-044.jpg}{Revisar nuevamente la receta de los cuatro componentes tras aprender mecanismos de aprendizaje}{Grabación oficial de Stanford, 00:49:00--00:49:30}

\begin{knowledgebox}{Un orden de decisiones ejecutable}
Primero audite los datos y las métricas de evaluación, luego use un proxy a pequeña escala para estimar calidad/repetición/límite de fase; posteriormente pruebe el front-loading con un presupuesto fijo; solo cuando la línea base, los FLOPs y el flujo post-entrenamiento estén estables, introduzca los rollouts de RLP. En cada paso conserve las opciones no adoptadas y los resultados fallidos, evitando convertir el ajuste de parámetros de la receta en alquimia irreproducible.
\end{knowledgebox}

\subsection{Preguntas del curso: detalles de implementación y límites}

Las preguntas complementan varios puntos clave. La maquinaria de ventaja de RLP es similar a GRPO; la principal variación es la recompensa de ganancia de información. Los experimentos seleccionan tokens aleatoriamente en el documento, no basándose en un cribado preciso por entropía. RLP usa una mezcla general de preentrenamiento, incluyendo gran cantidad de web crawl, y no opera solo sobre datos de razonamiento de alta calidad. También se observaron beneficios relativos a CPT desde checkpoints iniciales del orden del 20\% del preentrenamiento completo.

\teachervoice{00:49:54--00:50:35, el orador explica que RLP reutiliza la ventaja relativa grupal estilo GRPO; el foco de innovación es la construcción de recompensa.}

\teachervoice{00:55:36--00:57:43, el orador corrige al preguntante: RLP no funciona solo sobre datos de más alta calidad; en la práctica se seleccionan tokens aleatoriamente y también se realizaron experimentos favorables en checkpoints tempranos del orden de 4 T-token.}

\begin{lstlisting}[caption={Tabla mínima de experimento recomendada para evaluar RLP}]
experiment = {
    "start_checkpoint_tokens": ...,   # e.g. 19.8T
    "extra_data_tokens": ...,
    "rollouts_per_position": ...,
    "mean_thought_length": ...,
    "estimated_training_flops": ...,
    "gpu_hours": ...,
    "token_selection": "uniform_random",
    "data_mixture_version": ...,
    "post_training_recipe": ...,
    "benchmark_harness": ...,
}
\end{lstlisting}

\subsection{Límites de calidad de datos, RLHF y multimodal}

Sobre la calidad general de las páginas web, el orador menciona un clasificador educativo estilo FineWeb-Edu: usar una rúbrica para dar puntajes de 1--5 a los documentos, escalable además al dominio y nivel educativo. Respecto al sesgo de RLHF, alucinaciones y hacking de recompensa, ella indica claramente que es investigación continua de alineamiento, no una conclusión directa de las tres trabajos de esta conferencia. Para el alineamiento visión-lenguaje, su respuesta es que aún no se ha testeado.

\begin{warningbox}{No reescriba las respuestas cautelosas del Q\&A como conclusiones definitivas}
"RLHF aún puede ocupar una pequeña parte del post-entrenamiento", "el clasificador necesita actualización continua", "RLAIF podría reducir cierta subjetividad" son discusiones de clase y juicios empíricos, no validaciones experimentales de esta conferencia. Las apuntes mantienen sus límites, evitando escribir opiniones de áreas adyacentes como declaraciones de efectos de RLP.
\end{warningbox}

\subsection{Resumen de la sección}

Lo más importante desde el punto de vista de ingeniería no es repetir los headline gains, sino establecer un plano de control completo: calidad de datos y currículo, asignación de razonamiento, recompensa de rollout, estabilidad EMA, coincidencia token/FLOP/wall-clock, persistencia post-entrenamiento y restricciones transdominio. Solo con toda esta información disponible se puede juzgar si RLP o front-loading valen la pena ingresar al entrenamiento de producción.
\section{Síntesis y ampliación}

\subsection{Un hilo conductor que atraviesa tres líneas de trabajo}

El \textbf{preentrenamiento en dos fases} demuestra que el \textbf{orden} de los datos tiene valor; el \textbf{front-loading del razonamiento} muestra que el \textbf{momento} de los datos relacionados con la capacidad es relevante; y RLP indica que el \textbf{comportamiento y las recompensas} antes de la predicción son importantes. Juntos, estos tres enfoques impulsan el preentrenamiento desde un ajuste estático de máxima verosimilitud hacia un aprendizaje consciente del currículo, consciente de las etapas y consciente de la acción.

\begin{importantbox}{Marco de conocimiento final}
\begin{enumerate}
  \item La mezcla de datos debe registrar simultáneamente los pesos y el cronograma temporal.
  \item Calidad y época/repetición son variables distintas; incluso los datos de alta calidad tienen un límite de repetición.
  \item La asignación temprana de datos de razonamiento puede alterar el fundamento base disponible para la etapa posterior.
  \item RLP mide la ganancia informativa predictiva del pensamiento mediante contrafactuales sin pensar (no-think).
  \item Las recompensas densas, la línea base EMA y la actualización relativa al grupo determinan conjuntamente la estabilidad.
  \item Todos los beneficios deben interpretarse bajo las métricas de tokens, FLOPs, checkpoint, post-entrenamiento y evaluación en benchmarks.
\end{enumerate}
\end{importantbox}

\subsection{De los resultados del currículo a las preguntas de investigación}

En primer lugar, ¿es posible aprender un currículo continuo automáticamente en lugar de hacerlo manualmente en dos fases? En segundo lugar, ¿cómo se retroalimenta el sesgo del clasificador de calidad hacia el comportamiento del modelo? En tercer lugar, ¿es igualmente efectivo el front-loading para el razonamiento no STEM, multilingüe y multimodal? En cuarto lugar, ¿pueden incorporarse restricciones de facticidad, fidelidad causal o seguridad a la recompensa del pensamiento sin comprometer la escalabilidad? Y en quinto lugar, ¿puede reducirse el cómputo del rollout de RLP mediante pensamiento latente, generación especulativa o un batching más robusto?

\begin{knowledgebox}{Experimento factorial sugerido}
Es posible utilizar un diseño $2\times2\times2$ para activar y desactivar dos fases, front-loading y RLP por separado, midiendo los efectos principales y las interacciones bajo una arquitectura fija, un total de tokens constante, FLOPs aproximados y una receta de post-entrenamiento definida. Solo así se podrá responder si las tres estrategias son complementarias, redundantes o mutuamente interferentes, en lugar de ensamblar mecánicamente los resultados de tres artículos diferentes.
\end{knowledgebox}

\subsection{Lecturas adicionales}

Se recomienda leer siguiendo el orden de los problemas: primero, el preentrenamiento en dos fases de Feng et al., para comprender la mezcla de datos y el currículo; luego, el front-loading del razonamiento de Akter et al., verificando las condiciones de reason-base y las cinco lecciones; finalmente, RLP v2 de Hatamizadeh et al., centrándose en la identidad de mejora esperada, el algoritmo 1, el emparejamiento de cómputo y las ablaciones. Posteriormente, contrastar Quiet-STaR, RPT y RLPT comparando sus fuentes de recompensa y su granularidad.

\begin{warningbox}{Mantener la instantánea de la clase al leer los artículos}
Esta clase se impartió el 2026-04-30: dos fases utilizó arXiv v1, front-loading utilizó v1 y RLP utilizó v2 del 2026-03-01. Revisiones posteriores, reproducciones o resultados de producto pueden servir como nueva evidencia, pero no se deben retrotraer para convertirlos en hechos que ya estaban establecidos el día de la clase.
\end{warningbox}

Finalmente, el «inteligencia de próxima generación» propuesta por el curso no es un modelo nuevo aislado, sino una visión del entrenamiento: bajo datos y potencia computacional limitados, los modelos deben elegir con más inteligencia el orden de aprendizaje, establecer más tempranamente un fundamento de razonamiento y permitir que los comportamientos de exploración reciban retroalimentación escalable y comparable. El futuro del preentrenamiento podría no residir únicamente en un $N$ mayor, sino también en un proceso de aprendizaje mejorado.

\end{document}
